\section*{Резюме}

Настоящата дипломна работа разглежда проектирането, реализацията и оценяването на \ORMName{} библиотека за C++23~\cite{cpp23standard}. Основната идея на разработката е, че логическата структура на заявките и на модела на данните не трябва да се описва чрез низове, интерпретирани на етап изпълнение, а чрез типизирани \code{constexpr} обекти, валидирани от компилатора. По този начин част от традиционните грешки при достъпа до данни --- неправилни имена на полета, несъвместими типове, липсващи части от заявката и погрешно използване на неподдържани операции --- се преместват от фазата на изпълнение към фазата на компилация.

Разработената библиотека използва статичен модел на данните, изразен в C++ чрез конструкции като \code{property<T, Name>}, \code{relationship<...>}, \code{table_name_trait<T>} и типизиран query IR, представен чрез \code{select\_query}, \code{insert\_query}, \code{update\_query} и \code{delete\_query}. Този междинен модел е независим от конкретния тип база данни. Материализацията към релационни и нерелационни системи се осъществява чрез специализации на \code{connector\_trait<DB>}, които декларират поддържаните възможности и превеждат една и съща логическа заявка към SQL, BSON-подобни филтри, Redis команди, Cypher или CQL в зависимост от избрания backend.

В работата се анализират както релационните сценарии за SQLite, PostgreSQL и MySQL, така и нерелационните сценарии за MongoDB, Redis, Neo4j и Cassandra. Показано е как един и същ C++ модел диктува логическата схема, но се материализира различно като таблица, документ, ключ-стойност структура, графов възел или wide-column ред. Особено внимание е отделено на ролята на \code{store\_as::reference} и \code{store\_as::embed} като обща абстракция за връзки между обекти, която се интерпретира според модела на съхранение.

Верификацията на разработката се основава на unit и integration тестове от хранилището. Те покриват изграждането на query IR, capability-базираното ограничаване на операции, подготвените заявки, миграциите, нишковата безопасност и поведението на конкретните конектори при работа с реални или mock драйвери. Това позволява формулиране не само на архитектурен модел, но и на конкретни критерии за „напълно имплементиран и функциониращ“ конектор.

Основният принос на дипломната работа е в три направления. Първо, предложен е практически приложим модел за \ORMName{} библиотека, който съчетава статична типова безопасност и backend-независима архитектура. Второ, показано е как C++ описанието на данните може да се използва като първичен източник за логическа схема при различни типове бази данни. Трето, формализиран е connector contract, чрез който библиотеката може да се надгражда с нови backend реализации без промяна в user-facing query API.

\vspace{1cm}

\textbf{Ключови думи:} C++23, \ORMName{}, query IR, \code{connector\_trait}, типова безопасност, релационни и нерелационни бази данни, подготовени заявки, миграции
